\chapter{Lezione 8 - 17/10}
\section{L'intersezione di sottospazi è un sottospazio}
\Proposizione{L'intersezione di sottospazi è un sottospazio}{
    Sia \(V\) uno spazio vettoriale su un campo \(K\).
    Siano \(W_1, \dots, W_p \subseteq V\) sottospazi vettoriali di \(V\), con \(p \geq 2\).
    
    Definiamo \(W\) come la loro intersezione:
    \[ W = W_1 \cap \dots \cap W_p = \bigcap_{i=1}^p W_i \]
    
    Allora \(W\) è un sottospazio vettoriale di \(V\).

    \Dimostrazione{
        Per dimostrare che \(W\) è un sottospazio vettoriale, dobbiamo verificare le tre proprietà (il test dei sottospazi).

        \begin{enumerate}
            \item \textbf{Contiene il vettore nullo (\(W \neq \emptyset\)):}
            Poiché ogni \(W_i\) (per \(i=1, \dots, p\)) è un sottospazio, ognuno di essi contiene il vettore nullo:
            \[ \bar{0} \in W_i, \quad \forall i \in \{1, \dots, p\} \]
            Per definizione di intersezione, se \(\bar{0}\) appartiene a tutti gli insiemi, allora appartiene anche alla loro intersezione:
            \[ \bar{0} \in W = \bigcap_{i=1}^p W_i \]
            Questo dimostra che \(W\) non è l'insieme vuoto.
            
            \item \textbf{Chiusura rispetto alla somma:}
            Siano \(u, v \in W\). Dobbiamo dimostrare che \(u+v \in W\).
            Per definizione di \(W\), se \(u, v \in W\), allora \(u\) e \(v\) appartengono a ciascun \(W_i\):
            \[ u \in W_i \text{ e } v \in W_i, \quad \forall i \in \{1, \dots, p\} \]
            Poiché ogni \(W_i\) è un sottospazio, è chiuso rispetto alla somma. Dunque:
            \[ u+v \in W_i, \quad \forall i \in \{1, \dots, p\} \]
            Dato che \(u+v\) appartiene a tutti i sottospazi \(W_i\), appartiene anche alla loro intersezione:
            \[ u+v \in W \]
            
            
        \end{enumerate}
        
    }
}
\Proposizione{L'intersezione di sottospazi è un sottospazio}{
    \Dimostrazione{
        \begin{itemize}
            \item \textbf{Chiusura rispetto al prodotto per scalare:}
                Siano \(u \in W\) e \(\alpha \in K\). Dobbiamo dimostrare che \(\alpha u \in W\).
                Per definizione di \(W\), se \(u \in W\), allora \(u\) appartiene a ciascun \(W_i\):
                \[ u \in W_i, \quad \forall i \in \{1, \dots, p\} \]
                Poiché ogni \(W_i\) è un sottospazio, è chiuso rispetto al prodotto per scalare. Dunque:
                \[ \alpha u \in W_i, \quad \forall i \in \{1, \dots, p\} \]
                Dato che \(\alpha u\) appartiene a tutti i sottospazi \(W_i\), appartiene anche alla loro intersezione:
                \[ \alpha u \in W \]
        \end{itemize}
        Avendo verificato le tre proprietà, \(W\) è un sottospazio vettoriale di \(V\).
    }
}

\section{Sottospazio Somma}
\Definizione{Sottospazio Somma}{
    Siano \(W_1, \dots, W_p\) sottospazi vettoriali di \(V\).
    Si definisce \textbf{sottospazio somma} l'insieme:
    \[
    W_1 + \dots + W_p = \left\{ w_1 + \dots + w_p \mid w_i \in W_i \text{ per } i=1, \dots, p \right\}
    \]
    Questo è l'insieme di tutti i vettori di \(V\) che si possono scrivere come somma di vettori appartenenti a ciascun sottospazio.
}

\Proposizione{Il sottospazio somma è un sottospazio}{
    Siano \(W_1, \dots, W_p\) sottospazi vettoriali di \(V\).
    L'insieme \(W = W_1 + \dots + W_p\) è un sottospazio vettoriale di \(V\).
    
    \Dimostrazione{
        Dobbiamo verificare le tre proprietà dei sottospazi (test dei sottospazi).

        \begin{itemize}
            \item \textbf{Contiene il vettore nullo (\(W \neq \emptyset\)):}
            Poiché ogni \(W_i\) è un sottospazio, \(\bar{0} \in W_i\) per ogni \(i=1, \dots, p\).
            Possiamo quindi scrivere il vettore nullo come somma di elementi dei \(W_i\):
            \[ \bar{0} = \underbrace{\bar{0}}_{\in W_1} + \underbrace{\bar{0}}_{\in W_2} + \dots + \underbrace{\bar{0}}_{\in W_p} \]
            Dato che questa somma appartiene alla definizione di \(W\), si ha \(\bar{0} \in W\).
            Quindi \(W\) non è l'insieme vuoto.

            \item \textbf{Chiusura rispetto alla somma:}
            Siano \(u, v \in W\). Dobbiamo dimostrare che \(u+v \in W\).
            
            Per definizione di \(W\), se \(u, v \in W\), allora esistono:
            \[ u_1 \in W_1, \dots, u_p \in W_p \text{ tali che } u = u_1 + \dots + u_p \]
            \[ v_1 \in W_1, \dots, v_p \in W_p \text{ tali che } v = v_1 + \dots + v_p \]
            
            Calcoliamo la loro somma \(u+v\) e usiamo la proprietà associativa e commutativa della somma in \(V\) per riordinare i termini:
            \[ u+v = (u_1 + \dots + u_p) + (v_1 + \dots + v_p) = (u_1 + v_1) + \dots + (u_p + v_p) \]
            
            
            
        \end{itemize}
            
    }
}

\Proposizione{Il sottospazio somma è un sottospazio}{
    \Dimostrazione{
        Poiché ogni \(W_i\) è un sottospazio, è chiuso rispetto alla somma. Dunque:
            \[ (u_1 + v_1) \in W_1, \quad \dots, \quad (u_p + v_p) \in W_p \]
            
            Abbiamo scritto \(u+v\) come somma di elementi (dove l'i-esimo termine è \(u_i+v_i \in W_i\)).
            Per definizione, questo significa che \(u+v \in W\).
        \begin{itemize}

            

            \item \textbf{Chiusura rispetto al prodotto per scalare:}
            Siano \(u \in W\) e \(\lambda \in K\). Dobbiamo dimostrare che \(\lambda u \in W\).
            
            Come prima, \(u \in W \implies u = u_1 + \dots + u_p\), con \(u_i \in W_i\) per ogni \(i\).
            
            Calcoliamo il prodotto \(\lambda u\), usando la proprietà distributiva:
            \[ \lambda u = \lambda (u_1 + \dots + u_p) = (\lambda u_1) + \dots + (\lambda u_p) \]
            
            Poiché ogni \(W_i\) è un sottospazio, è chiuso rispetto al prodotto per scalare. Dunque:
            \[ (\lambda u_1) \in W_1, \quad \dots, \quad (\lambda u_p) \in W_p \]
            
            Abbiamo scritto \(\lambda u\) come somma di elementi (dove l'i-esimo termine è \(\lambda u_i \in W_i\)).
            Per definizione, questo significa che \(\lambda u \in W\).
        \end{itemize}
        Avendo verificato le tre proprietà, \(W\) è un sottospazio vettoriale di \(V\).
    }
}

\section{Somma di sottospazi generati da insiemi}
\Proposizione{Somma di sottospazi generati da insiemi}{
    Siano \( S_1, S_2, \dots, S_p \) insiemi tali che 
    \[
        W_1 = L(S_1), \; W_2 = L(S_2), \; \dots, \; W_p = L(S_p)
    \]
    siano sottospazi vettoriali di \( V \). Allora:
    \[
        W_1 + W_2 + \dots + W_p = L(S_1 \cup S_2 \cup \dots \cup S_p).
    \]
}
\Dimostrazione{
    \textbf{``\(\supseteq \)''} \\[2mm]
    Poiché
    \[
        S_1 \cup \dots \cup S_p \subseteq W_1 \cup \dots \cup W_p \subseteq W_1 + \dots + W_p,
    \]
    per definizione di chiusura lineare segue immediatamente che:
    \[
        L(S_1 \cup \dots \cup S_p) \subseteq W_1 + \dots + W_p.
    \]

    \medskip
}
\Dimostrazione{
    

    \textbf{``\(\subseteq \)''} \\[2mm]
    Sia 
    \[
        u \in W_1 + \dots + W_p.
    \]
    Allora esistono vettori 
    \[
        u_1 \in W_1, \; u_2 \in W_2, \; \dots, \; u_p \in W_p
    \]
    tali che
    \[
        u = u_1 + u_2 + \dots + u_p.
    \]

    Poiché ogni \( S_i \) genera \( W_i \), ciascun \( u_i \) può essere scritto come combinazione lineare di vettori di \( S_i \).  
    Cioè:

    \[
    \begin{aligned}
        &\exists\, v_1, \dots, v_t \in S_1, \; \alpha_1, \dots, \alpha_t \in K \quad \text{tali che} \quad u_1 = \alpha_1 v_1 + \dots + \alpha_t v_t, \\[1mm]
        &\exists\, w_1, \dots, w_s \in S_2, \; \beta_1, \dots, \beta_s \in K \quad \text{tali che} \quad u_2 = \beta_1 w_1 + \dots + \beta_s w_s, \\[1mm]
        &\dots \\[1mm]
        &\exists\, z_1, \dots, z_r \in S_p, \; \gamma_1, \dots, \gamma_r \in K \quad \text{tali che} \quad u_p = \gamma_1 z_1 + \dots + \gamma_r z_r.
    \end{aligned}
    \]

    Sommando queste combinazioni lineari otteniamo:
    \[
        u = \alpha_1 v_1 + \dots + \alpha_t v_t + \beta_1 w_1 + \dots + \beta_s w_s + \dots + \gamma_1 z_1 + \dots + \gamma_r z_r.
    \]

    Poiché tutti i vettori \( v_i, w_j, z_k \) appartengono all’unione \( S_1 \cup S_2 \cup \dots \cup S_p \),
    segue che \( u \in L(S_1 \cup S_2 \cup \dots \cup S_p) \).  
    Pertanto:
    \[
        W_1 + \dots + W_p \subseteq L(S_1 \cup S_2 \cup \dots \cup S_p).
    \]

    \medskip

    Combinando le due inclusioni otteniamo l’uguaglianza desiderata:
    \[
        W_1 + \dots + W_p = L(S_1 \cup S_2 \cup \dots \cup S_p).
    \]
}

\section{Regola di Grassman}

\Teorema{Regola di Grassmann}{
    Siano \(W_1\) e \(W_2\) sottospazi vettoriali di uno spazio \(V\) finitamente generato.
    Siano \(\dim W_1 = r\) e \(\dim W_2 = s\).
    
    Allora \(W_1 + W_2\) e \(W_1 \cap W_2\) sono anch'essi finitamente generati e vale la seguente relazione:
    \[
    \dim(W_1 + W_2) = \dim W_1 + \dim W_2 - \dim(W_1 \cap W_2)
    \]
    (ovvero \(\dim(W_1 + W_2) = r + s - i\), usando le notazioni della dimostrazione).

    \Dimostrazione{
    Sia \(W_{12} = W_1 \cap W_2\). Poiché \(W_{12}\) è un sottospazio di \(W_1\) (che è finitamente generato), anche \(W_{12}\) è finitamente generato.
    Sia \(\dim(W_{12}) = i\). Sia \(\mathcal{B}_{12} = \{u_1, \dots, u_i\}\) una sua base.

    Poiché \(W_{12} \subseteq W_1\), possiamo usare il \textbf{Teorema del Completamento della Base} per estendere \(\mathcal{B}_{12}\) a una base di \(W_1\).
    Sia \(\mathcal{B}_{W1} = \{u_1, \dots, u_i, v_{i+1}, \dots, v_r\}\) tale base (quindi \(\dim W_1 = r\)).

    Allo stesso modo, poiché \(W_{12} \subseteq W_2\), completiamo \(\mathcal{B}_{12}\) a una base di \(W_2\).
    Sia \(\mathcal{B}_{W2} = \{u_1, \dots, u_i, w_{i+1}, \dots, w_s\}\) tale base (quindi \(\dim W_2 = s\)).

    Ora consideriamo l'insieme \(\mathcal{B}_{SUM}\) dato dall'unione delle due basi:
    \[
    \mathcal{B}_{SUM} = \mathcal{B}_{W1} \cup \mathcal{B}_{W2} = \{u_1, \dots, u_i, v_{i+1}, \dots, v_r, w_{i+1}, \dots, w_s\}
    \]
    Il nostro obiettivo è dimostrare che \(\mathcal{B}_{SUM}\) è una base per \(W_1 + W_2\).
    
    \textbf{1. \(\mathcal{B}_{SUM}\) genera \(W_1 + W_2\)}
    Per definizione, il sottospazio somma \(W_1 + W_2\) è generato dall'unione dei sistemi di generatori:
    \[
    W_1 + W_2 = L(\mathcal{B}_{W1}) + L(\mathcal{B}_{W2}) = L(\mathcal{B}_{W1} \cup \mathcal{B}_{W2}) = L(\mathcal{B}_{SUM})
    \]
    Quindi \(\mathcal{B}_{SUM}\) è un sistema di generatori per \(W_1 + W_2\).

    \textbf{2. \(\mathcal{B}_{SUM}\) è linearmente indipendente}
    Questa è la parte cruciale. Dobbiamo dimostrare che l'unica combinazione lineare dei vettori di \(\mathcal{B}_{SUM}\) che dà il vettore nullo è quella con tutti i coefficienti nulli.
    Impostiamo l'equazione:
    \[
    \sum_{k=1}^i \alpha_k u_k + \sum_{j=i+1}^r \beta_j v_j + \sum_{l=i+1}^s \gamma_l w_l = \bar{0}
    \]
    Isoliamo i termini di \(\mathcal{B}_{W1}\) da un lato:
    \[
    \underbrace{\sum_{k=1}^i \alpha_k u_k + \sum_{j=i+1}^r \beta_j v_j}_{\text{chiamiamolo } z_1} = - \underbrace{\sum_{l=i+1}^s \gamma_l w_l}_{\text{chiamiamolo } z_2}
    \]
    Analizziamo il vettore \(z_1\): per costruzione, \(z_1\) è C.L. della base \(\mathcal{B}_{W1}\), quindi \(z_1 \in W_1\).
    
    Analizziamo il vettore \(z_2\): per costruzione, \(z_2\) è C.L. di vettori di \(\mathcal{B}_{W2}\), quindi \(z_2 \in W_2\).
    
    Dall'uguaglianza \(z_1 = -z_2\), deduciamo che il vettore \(z_1\) appartiene anche a \(W_2\) (essendo uguale a \(-z_2 \in W_2\)).
    
    Quindi \(z_1 \in W_1\) e \(z_1 \in W_2\), il che significa \(z_1 \in W_1 \cap W_2 = W_{12}\).
    
    Poiché \(z_1 \in W_{12}\), esso si deve poter scrivere come C.L. della base di \(W_{12}\), \(\mathcal{B}_{12} = \{u_1, \dots, u_i\}\).
    Esistono quindi \(\lambda_1, \dots, \lambda_i \in K\) tali che:
    \[
    z_1 = \sum_{k=1}^i \lambda_k u_k
    \]
    Ora confrontiamo le due espressioni che abbiamo per \(-z_2\):
    \[
    - \sum_{l=i+1}^s \gamma_l w_l = z_1 = \sum_{k=1}^i \lambda_k u_k
    \]
    Spostando tutto sullo stesso lato otteniamo:
    \[
    \bar{0} = \sum_{k=1}^i \lambda_k u_k + \sum_{l=i+1}^s \gamma_l w_l
    \]
    Questa è una combinazione lineare dei vettori di \(\mathcal{B}_{W2} = \{u_1, \dots, u_i, w_{i+1}, \dots, w_s\}\).
    Poiché \(\mathcal{B}_{W2}\) è una base, i suoi vettori sono linearmente indipendenti, e quindi tutti i coefficienti devono essere nulli:
    \[
    \lambda_1 = \dots = \lambda_i = 0 \quad \text{e} \quad \gamma_{i+1} = \dots = \gamma_s = 0
    \]
    
}

\Dimostrazione{
    Abbiamo dimostrato che tutti i \(\gamma_l\) sono nulli.
    
    Torniamo all'equazione di partenza. Sapendo che i \(\gamma_l\) sono nulli, essa si riduce a:
    \[
    \sum_{k=1}^i \alpha_k u_k + \sum_{j=i+1}^r \beta_j v_j = \bar{0}
    \]
    Questa è una C.L. dei vettori di \(\mathcal{B}_{W1} = \{u_1, \dots, u_i, v_{i+1}, \dots, v_r\}\).
    Poiché \(\mathcal{B}_{W1}\) è una base, anche i suoi vettori sono L.I., e quindi tutti i coefficienti devono essere nulli:
    \[
    \alpha_1 = \dots = \alpha_i = 0 \quad \text{e} \quad \beta_{i+1} = \dots = \beta_r = 0
    \]
    Avendo dimostrato che tutti i coefficienti \(\alpha, \beta, \gamma\) sono nulli, \(\mathcal{B}_{SUM}\) è linearmente indipendente.

    \textbf{3. Conclusione}
    Poiché \(\mathcal{B}_{SUM}\) è un sistema di generatori L.I. per \(W_1 + W_2\), è una sua base.
    La dimensione di \(W_1 + W_2\) è la cardinalità (il numero di vettori) di \(\mathcal{B}_{SUM}\):
    \[
    \dim(W_1 + W_2) = |\mathcal{B}_{SUM}| = \underbrace{i}_{(\text{vettori } u)} + \underbrace{(r-i)}_{(\text{vettori } v)} + \underbrace{(s-i)}_{(\text{vettori } w)}
    \]
    \[
    \dim(W_1 + W_2) = i + r - i + s - i = r + s - i
    \]
    Sostituendo i nomi delle dimensioni:
    \[
    \dim(W_1 + W_2) = \dim W_1 + \dim W_2 - \dim(W_1 \cap W_2)
    \]
}
}

\BoxBlue{Osservazione:}{Caso di Somma Diretta con intersezione banale}{
    Siano \(W_1, W_2\) sottospazi di \(V\).
    Se la loro intersezione è banale, cioè \(W_1 \cap W_2 = \{\bar{0}\}\), allora dalla Formula di Grassmann abbiamo:
    \[
    \dim(W_1 + W_2) = \dim W_1 + \dim W_2 - \dim(W_1 \cap W_2) = r + s - 0 = r + s
    \]
    Dalla dimostrazione di Grassmann, segue anche che l'unione delle basi \(\mathcal{B}_{W1}\) e \(\mathcal{B}_{W2}\) (che in questo caso non hanno vettori in comune, a parte lo \(\bar{0}\) che non è in nessuna base) forma una base per la somma:
    \[
    \mathcal{B}_{W_1 + W_2} = \mathcal{B}_{W_1} \cup \mathcal{B}_{W_2}
    \]
    La cardinalità della base è infatti \(|\mathcal{B}_{W_1 + W_2}| = |\mathcal{B}_{W_1}| + |\mathcal{B}_{W_2}| = r+s\).
}

\Definizione{Somma Diretta}{
    Siano \(W_1, \dots, W_p\) sottospazi vettoriali di \(V\), con \(p \geq 2\).
    
    La somma \(W = W_1 + \dots + W_p\) si dice \textbf{somma diretta} (e si scrive \(W = W_1 \oplus \dots \oplus W_p\)) se vale la seguente condizione:
    \[
    \forall i \in \{1, \dots, p\}, \quad W_i \cap \left( \sum_{j=1, j \neq i}^p W_j \right) = \{\bar{0}\}
    \]
    (Cioè, l'intersezione di ogni sottospazio con la somma di *tutti gli altri* è il solo vettore nullo).
}

\Proposizione{Caratterizzazione della Somma Diretta}{
    Sia \(W = W_1 + \dots + W_p\) un sottospazio somma (con \(p \geq 2\)).
    
    La somma è \textbf{diretta} (cioè \(W = W_1 \oplus \dots \oplus W_p\))
    \(\iff\)
    ogni vettore \(u \in W\) si scrive in \textbf{modo unico} come somma di vettori \(w_i \in W_i\).
    
    Ossia, se \(u \in W\) e abbiamo due decomposizioni:
    \begin{enumerate}
        \item \( u = u_1 + u_2 + \dots + u_p \quad \) (con \(u_i \in W_i\) per ogni \(i\))
        \item \( u = v_1 + v_2 + \dots + v_p \quad \) (con \(v_i \in W_i\) per ogni \(i\))
    \end{enumerate}
    allora le due decomposizioni devono essere identiche:
    \[
    u_1 = v_1, \quad u_2 = v_2, \quad \dots, \quad u_p = v_p
    \]
}

\Proposizione{Dimensione e Basi della Somma Diretta}{
    Siano \(W_1, \dots, W_p\) sottospazi vettoriali di \(V\), con \(\dim(W_i) = r_i\).
    
    Se la somma \(W = W_1 + \dots + W_p\) è \textbf{diretta} (cioè \(W = W_1 \boxplus \dots \boxplus W_p\)), e \(\mathcal{B}_1, \dots, \mathcal{B}_p\) sono basi di \(W_1, \dots, W_p\) rispettivamente, allora:
    \begin{enumerate}
        \item L'unione delle basi \(\mathcal{B} = \mathcal{B}_1 \cup \dots \cup \mathcal{B}_p\) è una base di \(W\).
        \item La dimensione della somma è la somma delle dimensioni:
        \[ \dim(W_1 \boxplus \dots \boxplus W_p) = r_1 + \dots + r_p \]
    \end{enumerate}
    
    \Dimostrazione{
        Dimostriamo per induzione sul numero \(p\) di sottospazi.

        \begin{itemize}
            \item \textbf{Base dell'induzione (p = 2):}
            Dobbiamo dimostrare la tesi per \(W = W_1 \boxplus W_2\).
            La somma è diretta, quindi per definizione \(W_1 \cap W_2 = \{\bar{0}\}\).
            Usando la Regola di Grassmann:
            \[ \dim(W_1 + W_2) = \dim(W_1) + \dim(W_2) - \dim(W_1 \cap W_2) \]
            \[ \dim(W) = r_1 + r_2 - 0 = r_1 + r_2 \]
            Come visto nell'Osservazione precedente, quando l'intersezione è banale, la base della somma è l'unione delle basi \(\mathcal{B} = \mathcal{B}_1 \cup \mathcal{B}_2\).
            La sua cardinalità è \(|\mathcal{B}| = |\mathcal{B}_1| + |\mathcal{B}_2| = r_1 + r_2\), che corrisponde alla dimensione.
            Il caso \(p=2\) è verificato.

            \item \textbf{Passo induttivo (p-1 \(\implies\) p):}
            Supponiamo (Ipotesi Induttiva) che la proposizione sia vera per \(p-1\) sottospazi.
            Cioè, per il sottospazio \(U = W_1 \boxplus \dots \boxplus W_{p-1}\):
            \begin{itemize}
                \item La sua base è \(\mathcal{B}_U = \mathcal{B}_1 \cup \dots \cup \mathcal{B}_{p-1}\).
                \item La sua dimensione è \(\dim(U) = r_1 + \dots + r_{p-1}\).
            \end{itemize}
            
            Ora consideriamo la somma a \(p\) termini:
            \[ W = W_1 \boxplus \dots \boxplus W_p = (W_1 \boxplus \dots \boxplus W_{p-1}) + W_p = U + W_p \]
            
            Per la definizione di somma diretta, l'intersezione di \(W_p\) con la somma di tutti gli altri è banale:
            \[ W_p \cap (W_1 + \dots + W_{p-1}) = W_p \cap U = \{\bar{0}\} \]
            
            Ci siamo ricondotti al caso base (somma di *due* sottospazi, \(U\) e \(W_p\), con intersezione nulla).
        \end{itemize}
    }
}

\Proposizione{Dimensione e Basi della Somma Diretta}{
    \Dimostrazione{
            \textbf{(Dimensione)} Applichiamo di nuovo Grassmann a \(U + W_p\):
            \[ \dim(W) = \dim(U + W_p) = \dim(U) + \dim(W_p) - \dim(U \cap W_p) \]
            Sostituendo i valori (dall'ipotesi induttiva e dai dati):
            \[ \dim(W) = (r_1 + \dots + r_{p-1}) + r_p - 0 = r_1 + \dots + r_p \]
            (La parte 2 della tesi è dimostrata).
            
            \textbf{(Base)} Sempre per il caso \(p=2\), la base della somma \(U \boxplus W_p\) è l'unione delle loro basi:
            \[ \mathcal{B} = \mathcal{B}_U \cup \mathcal{B}_p \]
            Sostituendo la base \(\mathcal{B}_U\) (dall'ipotesi induttiva):
            \[ \mathcal{B} = (\mathcal{B}_1 \cup \dots \cup \mathcal{B}_{p-1}) \cup \mathcal{B}_p = \mathcal{B}_1 \cup \dots \cup \mathcal{B}_p \]
            (La parte 1 della tesi è dimostrata).
            
            Avendo dimostrato il passo induttivo, la proposizione è vera per ogni \(p \geq 2\).
    }
}